% Content imported from Report_newest.docx.md with approved edits.
\chapter{KỸ THUẬT VÀ NGUYÊN LÝ KIỂM SOÁT TRÊN MÔ HÌNH POOL}\label{v4-chapter-3}

Thay vì tách rời việc mô tả từng tầng với các cơ chế bảo vệ, Chương~\ref{v4-chapter-3} đi theo bốn cụm tầng của mô hình Pool. Tại mỗi cụm tầng, phần phân tích chỉ ra cách vận hành khi dùng chung tài nguyên, giải thích nguyên nhân phát sinh rủi ro và trình bày các giải pháp kỹ thuật cho bốn khía cạnh K1--K4.

Thứ tự bốn cụm tầng trong Chương~\ref{v4-chapter-3} không liệt kê theo số thứ tự từ Tầng 1 đến Tầng 9 mà đi theo đúng đường đi của một yêu cầu trong hệ thống. Sau khi giới hạn phạm vi hạ tầng nền ở Tầng 1--3 (Mục~\ref{v4-sec-3-1}), chương đi vào cổng kiểm tra đầu tiên ở Tầng 7 (Mục~\ref{v4-sec-3-2}), tiếp nối vào môi trường chạy ứng dụng ở Tầng 8--9 (Mục~\ref{v4-sec-3-3}), và kết thúc ở lớp lưu trữ cuối cùng tại Cụm Cơ sở dữ liệu Tầng 4--6 (Mục~\ref{v4-sec-3-4}).

\section{Cụm hạ tầng nền tảng (Tầng 1-3)}\label{v4-sec-3-1}

Trong mô hình Pool Cấp độ 3, toàn bộ Cụm Hạ tầng nền (Tầng 1--3) được tổ chức theo hướng gộp dùng chung. Mọi thuê bao chia sẻ chung máy chủ, mạng nội bộ và nền tảng vùng chứa. Chia sẻ máy chủ có thể tận dụng năng lực nhàn rỗi tốt hơn cấp phát riêng cho từng thuê bao nhưng tạo ra cạnh tranh tài nguyên vật lý giữa các dịch vụ trên cùng máy chủ.

Cần phân biệt rõ giữa cô lập hạ tầng không biết thuê bao ở Tầng 1--3 và cô lập phần mềm có nhận biết thuê bao ở Tầng 4--9. Các công cụ hệ điều hành như Linux cgroups hay Namespaces chỉ giới hạn CPU và RAM ở mức từng vùng chứa (Container/Pod), không phân biệt được thuê bao nào đang chạy bên trong cùng một tiến trình ứng dụng hay cơ sở dữ liệu.

Do đó, ở Tầng 1--3, đề tài chỉ dùng công cụ điều phối vùng chứa để đóng gói dịch vụ và đặt trần tài nguyên CPU, RAM qua Docker Compose (cgroups). Trọng tâm giải pháp đa thuê bao nằm ở các cơ chế phần mềm từ Tầng 4 đến Tầng 9.

\section{Cụm phần mềm trung gian (Tầng 7)}\label{v4-sec-3-2}

Nằm giữa lớp ứng dụng (Tầng 8--9) và cơ sở dữ liệu (Tầng 4--6) là Tầng 7 -- Phần mềm trung gian (Middleware). Tầng này gồm cổng biên (API Gateway / Reverse Proxy), bộ nhớ đệm phân tán, hàng đợi thông điệp, kênh phát tán sự kiện và kho lưu trữ tệp. Khi áp dụng mô hình Pool tại Tầng 7, toàn bộ lưu lượng của mọi tổ chức đều đi qua một đường ống biên dùng chung như minh họa trên Hình~\ref{fig:v4-3-1}:

\VfourFigure{3.1}{assets/figures/v4/image17.png}{Tổ chức mô hình Pool tại Tầng 7 (Middleware).}{fig:v4-3-1}

Vì Tầng 7 vừa là cửa ngõ tiếp nhận toàn bộ yêu cầu vừa giữ bộ nhớ đệm tốc độ cao dùng chung, việc gộp tài nguyên tại đây làm nảy sinh đồng thời ba vấn đề kỹ thuật:

\begin{itemize}
\item
  \textbf{Thứ nhất, nguy cơ giả mạo mã thuê bao ở cổng biên (K1):} Khi mọi tổ chức gọi vào cùng một cổng tiếp nhận chung, máy khách có thể tự gửi mã thuê bao qua đường dẫn URL, phần thân JSON hoặc tiêu đề HTTP thông thường. Nếu máy chủ tin ngay giá trị tự khai này mà không đối chiếu với thẻ bài có chữ ký số, người dùng của tổ chức này có thể đổi tham số để xem hoặc sửa dữ liệu của tổ chức khác (lỗ hổng IDOR).
\item
  \textbf{Thứ hai, nguy cơ đọc nhầm khóa đệm và cạn bộ nhớ đệm do hàng xóm ồn ào (K2 \& K3):} Nếu các khóa trong bộ nhớ đệm dùng chung chỉ đặt theo mã bản ghi (như \textless entity\textgreater:\textless id\textgreater) mà thiếu tiền tố tổ chức, dữ liệu giữa các thuê bao sẽ bị ghi đè hoặc đọc chéo (K2). Về mặt hiệu năng (K3), khi một thuê bao gửi lưu lượng tăng đột biến hoặc quét lượng dữ liệu lịch sử lớn, yêu cầu đó sẽ chiếm băng thông và đẩy dữ liệu hay dùng của các thuê bao khác ra khỏi bộ nhớ đệm.
\item
  \textbf{Thứ ba, nguy cơ lệch cấu hình giữa các tiến trình ứng dụng (K4):} Khi quản trị viên một tổ chức thay đổi cấu hình trong lúc hệ thống đang chạy, các bản sao lưu sẵn trên từng tiến trình ở Tầng 8 sẽ bị cũ đi nếu không nhận được tín hiệu xóa bỏ kịp thời qua Tầng 7, dẫn đến xử lý sai quy tắc mới.
\end{itemize}

Nếu chọn mô hình Silo hay Bridge ở Tầng 7 để tách riêng không gian đệm và hàng đợi cho từng thuê bao, chi phí quản trị sẽ tăng cao khi số lượng khách hàng lớn, đồng thời nhiều hệ thống đệm hay hàng đợi cũng giới hạn số lượng phân vùng logic (như số cơ sở dữ liệu logic trong Redis hay số Virtual Host). Vì vậy, mô hình Pool giải quyết các rủi ro tại Tầng 7 bằng những giải pháp phần mềm trình bày trong Mục~\ref{v4-sec-3-2-1} và Mục~\ref{v4-sec-3-2-2}.

\subsection{Xác thực định danh mật mã tại biên}\label{v4-sec-3-2-1}

Ở cửa ngõ Tầng 7, hệ thống không tin vào mã thuê bao do trình duyệt hay máy khách tự khai trong URL hoặc tiêu đề HTTP thông thường. Nguồn thông tin đáng tin cậy duy nhất để xác định thuê bao là thẻ bài định danh được ký số mật mã ở phía máy chủ.

Cụ thể, trước khi đi vào Tầng 8, mọi yêu cầu đều phải qua ba bước kiểm tra tại bộ lọc biên:

\begin{enumerate}
\def\labelenumi{\arabic{enumi}.}
\item
  \textbf{Kiểm tra chữ ký số và tính toàn vẹn của thẻ bài:} Khi người dùng đăng nhập vào không gian làm việc của một tổ chức, máy chủ cấp thẻ bài có chữ ký số (ví dụ JWT) chứa mã người dùng, mã thuê bao đang hoạt động (tenant\_\allowbreak{}id) và vai trò tương ứng. Ở các yêu cầu sau đó, bộ lọc biên kiểm tra chữ ký và hạn dùng của thẻ bài; mọi sửa đổi tenant\_\allowbreak{}id ở phía máy khách đều làm sai chữ ký và bị từ chối ngay (trả mã HTTP 401).
\item
  \textbf{Đối chiếu chéo tên miền phụ và tư cách thành viên:} Nếu yêu cầu có gửi kèm tên miền phụ (Host) hoặc tiêu đề phụ trợ (như X-Tenant-ID), bộ lọc biên chỉ xem đó là tín hiệu gợi ý định hướng. Hệ thống phải đối chiếu giá trị gợi ý này với tenant\_\allowbreak{}id đã ký trong thẻ bài, đồng thời xác nhận tài khoản vẫn đang là thành viên hợp lệ của tổ chức đó.
\item
  \textbf{Tạo ngữ cảnh tin cậy chuyển xuống Tầng 8:} Chỉ sau khi vượt qua hai bước kiểm tra trên, mã tenant\_\allowbreak{}id mới được nạp vào ngữ cảnh xử lý của Tầng 8 (Mục~\ref{v4-sec-3-3-1}). Cách làm này chặn ngay các lỗi giả mạo tham chiếu IDOR ngay tại cổng vào.
\end{enumerate}

\subsection{Điều tiết lưu lượng, bộ nhớ đệm hai tầng và đồng bộ sự kiện}\label{v4-sec-3-2-2}

Để ngăn hiện tượng hàng xóm ồn ào (K3), tách riêng không gian đệm theo từng khách hàng (K2) và giữ cấu hình đồng bộ giữa các tiến trình (K4), Tầng 7 phối hợp cùng Tầng 8 qua ba cơ chế chính:

Thứ nhất, giới hạn tốc độ truy cập riêng cho từng thuê bao tại biên (K3). Thay vì đặt chung một hạn mức lưu lượng cho cả máy chủ (vốn vẫn để mặc một khách hàng chiếm hết băng thông), Tầng 7 áp dụng bộ đếm hạn mức độc lập theo từng tenant\_\allowbreak{}id trên bộ nhớ tốc độ cao bằng thuật toán Thùng thẻ bài (Token Bucket) hoặc Cửa sổ trượt (Sliding Window). Mỗi thuê bao được cấp hạn mức riêng theo gói cước và bị trừ dần trên từng yêu cầu. Khi một thuê bao gửi lượng yêu cầu vượt quá ngưỡng, hệ thống từ chối ngay phần dư thừa tại Tầng 7 (trả mã HTTP 429) để giữ đường ống thông suốt cho các thuê bao còn lại.

Thứ hai, đặt tiền tố không gian tên cho khóa đệm (K2) kết hợp bộ nhớ đệm hai tầng dùng thuật toán W-TinyLFU (K3). Để vừa truy xuất nhanh vừa dùng chung dữ liệu giữa nhiều tiến trình, hệ thống kết hợp bộ nhớ đệm cục bộ L1 ngay trong tiến trình Tầng 8 với bộ nhớ đệm phân tán L2 ở Tầng 7. Mọi khóa lưu trữ đều bắt buộc phải có tiền tố dạng tenant:\{tenant\_\allowbreak{}id\}:\{domain\_\allowbreak{}entity\}:\{entity\_\allowbreak{}id\} nhằm tránh đọc nhầm dữ liệu giữa các tổ chức.

Tuy nhiên, thuật toán thay thế LRU (Least Recently Used) truyền thống chỉ xét thời điểm truy cập gần nhất nên dễ bị đánh bật dữ liệu nóng khi một thuê bao bất ngờ quét khối lượng dữ liệu lớn. Để khắc phục tình trạng này, hệ thống sử dụng thuật toán Window TinyLFU (W-TinyLFU) \cite{v4-20} kết hợp cả độ mới lẫn tần suất truy cập lịch sử như minh họa trên Hình~\ref{fig:v4-3-2}:

\VfourFigure{3.2}{assets/figures/v4/image18.png}{Kiến trúc bộ nhớ đệm hai tầng (L1/L2) với W-TinyLFU và đồng bộ Pub/Sub.}{fig:v4-3-2}

Trên Hình~\ref{fig:v4-3-2}, W-TinyLFU đặt một cửa sổ đón nhỏ (Window Cache) trước vùng lưu chính và dùng bộ lọc tần suất để quyết định bản ghi nào đủ điều kiện ở lại. Nhờ vậy, dù một thuê bao đọc lướt qua nhiều dữ liệu ít dùng, các cấu hình được truy cập thường xuyên của các thuê bao khác vẫn giữ nguyên trong bộ nhớ đệm.

Lựa chọn bộ đệm hai tầng L1 (W-TinyLFU) kết hợp L2 phân tán được ưu tiên thay cho một tầng đệm LRU thông thường vì hai lý do thực tế: đọc siêu dữ liệu cấu hình từ L1 giúp bỏ qua trễ mạng nội bộ, còn W-TinyLFU giúp bộ đệm không bị quét sạch bởi các truy vấn bất thường.

Thứ ba, đồng bộ xóa bộ đệm qua kênh phát sự kiện Pub/Sub (K4). Vì bộ đệm L1 nằm riêng trong bộ nhớ của từng tiến trình ứng dụng ở Tầng 8, khi một thuê bao sửa cấu hình, tiến trình xử lý lệnh sửa sẽ xóa khóa tương ứng trên L2 rồi gửi một thông báo nhẹ mang mã tenant\_\allowbreak{}id lên kênh Publish/Subscribe ở Tầng 7. Mọi tiến trình đang lắng nghe kênh này lập tức xóa bản sao cấu hình cũ của đúng tenant\_\allowbreak{}id đó trong L1, bảo đảm cấu hình mới có hiệu lực đồng nhất ngay lúc chạy.

\section{Cụm ứng dụng và môi trường thực thi (Tầng 8-9)}\label{v4-sec-3-3}

Cụm Ứng dụng và Môi trường thực thi gồm Tầng 8 (máy chủ ứng dụng) và Tầng 9 (thể hiện ứng dụng đang xử lý nghiệp vụ). Trong nguyên mẫu Pool cấp độ 3, các thuê bao dùng chung bộ mã nguồn triển khai và thể hiện ứng dụng tại Tầng 9. Các yêu cầu chạy xen kẽ trên môi trường thực thi và bể luồng dùng chung tại Tầng 8, như minh họa ở Hình~\ref{fig:v4-3-3}:

\VfourFigure{3.3}{assets/figures/v4/image19.png}{Tổ chức mô hình Pool tại Tầng 8-9 (Ứng dụng và môi trường thực thi).}{fig:v4-3-3}

Việc chạy chung nhiều thuê bao trong cùng một tiến trình và một bộ mã nguồn đặt ra ba vấn đề kỹ thuật:

\begin{itemize}
\item
  \textbf{Nguy cơ lẫn ngữ cảnh giữa các luồng xử lý (K1 -- Thách thức 1):} Để xử lý nhanh, máy chủ ứng dụng giữ sẵn một Bể luồng (Thread Pool) tái sử dụng cho nhiều yêu cầu nối tiếp nhau. Nếu một luồng vừa phục vụ Thuê bao A mà không xóa sạch mã thuê bao cũ khi kết thúc, yêu cầu tiếp theo của Thuê bao B dùng lại luồng đó sẽ mang nhầm ngữ cảnh của Thuê bao A.
\item
  \textbf{Khó khăn khi tùy biến trên một bộ mã nguồn chung (K4 -- Thách thức 4):} Hệ thống không được tách nhánh mã nguồn hay viết câu lệnh rẽ nhánh cứng theo từng tên khách hàng, nhưng vẫn phải cho phép mỗi tổ chức dùng giao diện, luật kiểm tra và trường dữ liệu khác nhau ngay khi đang chạy.
\item
  \textbf{Nguy cơ tràn bộ nhớ chung và nghẽn thu gom rác (K3 -- Thách thức 3):} Mọi luồng xử lý cùng cấp phát đối tượng trên một vùng nhớ Heap dùng chung ở Tầng 8. Nếu một thuê bao truy vấn danh sách dữ liệu quá lớn không có phân trang, bộ nhớ Heap nhanh chóng đầy lên và kích hoạt thu gom rác toàn cục, làm khựng lại mọi yêu cầu của các thuê bao láng giềng.
\end{itemize}

Nếu tách tiến trình chạy riêng cho mỗi thuê bao theo mô hình Silo hay Bridge (chạy Worker hoặc Pod độc lập cho từng khách hàng), vùng nhớ Heap và bể luồng được cô lập ở cấp hệ điều hành. Nhưng mỗi tiến trình riêng đều tốn một lượng RAM nền đáng kể, làm tăng nhu cầu bộ nhớ nền theo số tiến trình chạy đồng thời.

\subsection{Quản lý ngữ cảnh theo phạm vi thực thi}\label{v4-sec-3-3-1}

Trong phần mềm nhiều tầng, việc truyền thủ công biến tenant\_\allowbreak{}id qua từng hàm sẽ làm mã nguồn cồng kềnh và dễ truyền nhầm tham số; ngược lại, đặt vào biến dùng chung toàn cục sẽ bị các luồng ghi đè lên nhau. Do đó, Tầng 8 lưu mã thuê bao vào ngữ cảnh riêng gắn theo từng luồng hoặc chuỗi thực thi của yêu cầu (K1).

Cụ thể, hệ thống dùng cấu trúc lưu ngữ cảnh cục bộ theo phạm vi yêu cầu (như ThreadLocal trong Java/Spring Boot, AsyncLocalStorage trong Node.js hay context.Context trong Go). Khi yêu cầu vượt qua bộ lọc biên, mã thuê bao đã xác thực được gắn chặt vào ngữ cảnh riêng của luồng xử lý đó suốt thời gian chạy.

Để tránh sót dữ liệu khi Bể luồng tái sử dụng luồng cũ cho khách hàng khác, cơ chế quản lý ngữ cảnh tuân theo ba quy tắc bắt buộc. Thứ nhất, đối tượng ngữ cảnh (chứa tenant\_\allowbreak{}id, user\_\allowbreak{}id, vai trò) được tạo ở dạng chỉ đọc để logic nghiệp vụ ở Tầng 9 không thể tự ý sửa đổi giữa chừng. Thứ hai, bộ lọc đầu vào phải bọc toàn bộ lệnh gọi phía sau trong khối try ... finally để luôn xóa sạch ngữ cảnh trước khi trả luồng về Bể luồng, kể cả khi phát sinh lỗi ngoại lệ. Thứ ba, khi tách sang luồng chạy nền, cần sao chép ngữ cảnh từ luồng cha sang luồng con lúc khởi tạo và xóa sạch sau khi chạy xong (mẫu Context Decorator). Toàn bộ vòng đời này được minh họa trên Hình~\ref{fig:v4-3-4}:

\VfourFigure{3.4}{assets/figures/v4/image20.png}{Biểu đồ tuần tự quản lý và thu hồi ngữ cảnh luồng tại Tầng 8.}{fig:v4-3-4}

Cách gắn ngữ cảnh theo luồng kèm khối bảo vệ finally được chọn thay cho việc truyền biến tenant\_\allowbreak{}id qua toàn bộ chữ ký hàm nhằm giữ mã nguồn nghiệp vụ gọn sạch, giảm sai sót do truyền nhầm tham số mà vẫn khóa được rủi ro rò rỉ khi tái sử dụng luồng.

\subsection{Đánh giá quy tắc và tùy biến động theo siêu dữ liệu}\label{v4-sec-3-3-2}

Để thay đổi giao diện và cách xử lý cho từng khách hàng trên một bộ mã nguồn chung ở Tầng 9, hệ thống tách phần mã thực thi cố định khỏi phần quy tắc cấu hình của từng thuê bao. Mã nguồn chỉ cài đặt khung xử lý tổng quát, còn các quy định riêng của mỗi tổ chức được lưu thành siêu dữ liệu (Metadata) để đọc và áp dụng trực tiếp lúc chạy.

Thay vì viết các điều kiện kiểm tra cố định trong mã lệnh, mỗi điểm thay đổi nghiệp vụ được biểu diễn thành một hàm tính hành vi theo ngữ cảnh hiện tại:

\textbf{\emph{Behavior = f(PolicyKey, EvaluationContext, MetadataRuleset)}}

Trong biểu thức trên, PolicyKey là tên quy tắc cần tra cứu (như chế độ giới hạn công việc hay cấu hình giao diện), EvaluationContext là thông tin thuê bao hiện tại lấy từ TenantContext (tenant\_\allowbreak{}id, gói cước, vai trò), còn MetadataRuleset là kho cấu hình siêu dữ liệu. Động cơ Đánh giá Quy tắc (Rule Evaluation Engine) tra cứu theo ba mức ưu tiên từ hẹp đến rộng:

\begin{enumerate}
\def\labelenumi{\arabic{enumi}.}
\item
  \textbf{Mức 1 (Ưu tiên cấu hình riêng của thuê bao):} Kiểm tra xem tenant\_\allowbreak{}id hiện tại có thiết lập giá trị riêng cho PolicyKey đó hay không.
\item
  \textbf{Mức 2 (Theo gói dịch vụ):} Nếu thuê bao không cấu hình riêng, áp dụng quy tắc chung của gói cước mà tổ chức đang sử dụng.
\item
  \textbf{Mức 3 (Mặc định hệ thống):} Dùng cấu hình mặc định an toàn nếu hai mức trên không có thiết lập riêng.
\end{enumerate}

Nhờ động cơ đánh giá siêu dữ liệu này, ứng dụng SaaS ở Tầng 9 tùy biến được đồng thời trên ba lớp: (1) giao diện hiển thị, (2) quy trình và luật nghiệp vụ, và (3) cấu trúc trường dữ liệu mở rộng, như thể hiện trên Hình~\ref{fig:v4-3-5}:

\VfourFigure{3.5}{assets/figures/v4/image21.png}{Nguyên lý động cơ đánh giá quy tắc và tùy biến động ba lớp.}{fig:v4-3-5}

\subsection{Bảo vệ vùng nhớ và phân bổ luồng thực thi}\label{v4-sec-3-3-3}

Việc chia sẻ vùng nhớ heap và bể luồng tại Tầng 8--9 đặt ra nhu cầu khống chế kích thước yêu cầu, thời gian chiếm luồng và tác động của bộ thu gom rác. Ba biện pháp sau là nguyên tắc thiết kế. Chương~\ref{v4-chapter-4} mô tả cơ chế đã hiện thực và Chương~\ref{v4-chapter-5} báo cáo kết quả đo; thực nghiệm hiện có chưa đánh giá riêng giới hạn luồng và bộ nhớ theo từng thuê bao.

Về mặt nguyên lý thiết kế, việc bảo vệ tài nguyên chạy ở Tầng 8--9 dựa trên ba biện pháp chính:

\begin{enumerate}
\def\labelenumi{\arabic{enumi}.}
\item
  \textbf{Bắt buộc phân trang và giới hạn kích thước dữ liệu:} Mọi API trả về danh sách đều phải áp dụng phân trang và giới hạn số bản ghi tối đa cho mỗi lần gọi. Cách này tránh việc một thuê bao nạp đồng thời lượng bản ghi quá lớn vào vùng nhớ Heap.
\item
  \textbf{Giới hạn thời gian chạy và số luồng song song theo thuê bao:} Mỗi yêu cầu được gán ngưỡng thời gian tối đa để tự động ngắt khi kẹt tài nguyên, kết hợp đặt giới hạn số yêu cầu được phép chạy cùng lúc cho từng tenant\_\allowbreak{}id. Khi một thuê bao gửi quá nhiều tác vụ chậm cùng lúc, phần vượt ngưỡng sẽ bị từ chối sớm để chừa luồng phục vụ các thuê bao khác.
\item
  \textbf{Đẩy tác vụ nặng sang xử lý nền bất đồng bộ:} Các thao tác tốn nhiều CPU và bộ nhớ như xuất báo cáo lớn hoặc nhập dữ liệu hàng loạt được tách khỏi luồng xử lý tương tác chính sang hàng đợi chạy nền có giới hạn riêng, vẫn gắn kèm ngữ cảnh thuê bao an toàn theo Mục~\ref{v4-sec-3-3-1}.
\end{enumerate}

\section{Cụm cơ sở dữ liệu (Tầng 4-6)}\label{v4-sec-3-4}

Phần lưu trữ dữ liệu có trạng thái của toàn hệ thống đặt tại Cụm Cơ sở dữ liệu (Tầng 4--6). Theo mô hình Pool đã chọn, mọi thuê bao dùng chung một cụm máy chủ hệ quản trị CSDL (Tầng 4), một cơ sở dữ liệu logic (Tầng 5) và một bộ bảng vật lý ở Tầng 6. Dữ liệu của các khách hàng nằm xen kẽ trong cùng một bảng và được phân biệt với nhau nhờ cột định danh tenant\_\allowbreak{}id như trên Hình~\ref{fig:v4-3-6}:

\VfourFigure{3.6}{assets/figures/v4/image22.png}{Tổ chức mô hình Pool và bảo mật cấp dòng tại Tầng 4-6.}{fig:v4-3-6}

Việc gom dữ liệu về một tập bảng chung như trên Hình~\ref{fig:v4-3-6} giúp tái sử dụng một bể kết nối duy nhất và quản lý di trú tập trung trên một lược đồ. Thời gian chạy di trú còn phụ thuộc dữ liệu và thao tác thực hiện. Tuy vậy, gộp chung lưu trữ cũng làm bộc lộ bốn rủi ro kỹ thuật lớn xuyên suốt K1--K4:

\begin{itemize}
\item
  \textbf{Rủi ro giữ nhầm ngữ cảnh trên bể kết nối CSDL dùng chung (K1):} Tầng 8 dùng chung một Bể kết nối vật lý (Connection Pool) giữ kết nối mở sẵn tới Tầng 4 để tái sử dụng cho hàng loạt giao dịch nối tiếp nhau. Nếu mã thuê bao gán vào kết nối nhưng không gắn với vòng đời giao dịch để tự xóa khi xong, giao dịch sau lấy lại kết nối đó sẽ đọc hoặc ghi nhầm vào dữ liệu của thuê bao trước.
\item
  \textbf{Rủi ro lộ dữ liệu do lập trình viên viết sót điều kiện lọc (K2):} Khi bản ghi của mọi tổ chức nằm cùng một bảng vật lý, nếu hệ thống chỉ dựa vào việc gắn thủ công điều kiện WHERE tenant\_\allowbreak{}id = ? trong mã nguồn, một câu truy vấn phức tạp bị viết thiếu điều kiện lọc có thể làm lộ dữ liệu của tổ chức này cho tổ chức khác.
\item
  \textbf{Xung đột giữa cấu trúc bảng cố định và nhu cầu thêm trường động (K4):} Trên bảng dùng chung ở Tầng 6, hệ thống không thể chạy lệnh DDL (như thêm cột vật lý mới) cho từng khách hàng lúc đang chạy vì thao tác này sẽ khóa bảng trên toàn hệ thống và sinh ra nhiều cột trống lãng phí.
\item
  \textbf{Suy giảm tốc độ truy vấn và nghẽn I/O khi bảng lớn lên (K3):} Khi dữ liệu nhiều khách hàng dồn vào cùng một bảng, nếu thiếu chỉ mục đặt cột tenant\_\allowbreak{}id ở đầu, truy vấn của một thuê bao có thể phải quét cả bảng chung, làm tốn băng thông đọc đĩa và đẩy trang dữ liệu của các thuê bao khác khỏi bộ nhớ đệm CSDL.
\end{itemize}

Để làm rõ lý do chọn Pool thay cho các mô hình khác tại Tầng 4--6, Hình~\ref{fig:v4-3-7} và Hình~\ref{fig:v4-3-8} minh họa hai phương án đối chứng. Ở mô hình Silo (Hình~\ref{fig:v4-3-7} -- tách máy chủ hệ quản trị CSDL hoặc cơ sở dữ liệu riêng), ranh giới vật lý rất rõ nhưng mỗi cơ sở dữ liệu phải giữ một bể kết nối riêng làm lượng kết nối và RAM tăng theo O(N). Ở mô hình Bridge (Hình~\ref{fig:v4-3-8} -- tách lược đồ riêng cho từng thuê bao kết hợp lệnh chuyển không gian tên như search\_\allowbreak{}path), các thuê bao dùng chung được bể kết nối; nhưng khi mở rộng lên hàng nghìn thuê bao, số lượng bảng vật lý tăng vọt làm phình bảng danh mục hệ thống (như pg\_\allowbreak{}class) và việc nâng cấp bảng vẫn phải lặp O(N) \cite{v4-21}. Từ so sánh đó, Mục~\ref{v4-sec-3-4-1} và 3.4.2 trình bày giải pháp bảo vệ trực tiếp cho mô hình Pool tại Tầng 4--6.

\VfourFigure{3.7}{assets/figures/v4/image23.png}{Liên hệ mô hình Silo tại Tầng 4-5.}{fig:v4-3-7}

\VfourFigure{3.8}{assets/figures/v4/image24.png}{Liên hệ mô hình Bridge (tách lược đồ) tại Tầng 6.}{fig:v4-3-8}

\subsection{Bảo mật cấp dòng (RLS) và biến phiên giao dịch}\label{v4-sec-3-4-1}

Khi lưu chung nhiều khách hàng trong cùng một bảng, hệ thống không thể chỉ trông chờ vào câu lệnh lọc WHERE ở tầng ứng dụng vì lập trình viên vẫn có thể viết sót. Vì vậy, Tầng 6 thiết lập rào chắn bắt buộc ngay trong nhân hệ quản trị CSDL bằng cơ chế Bảo mật cấp dòng (Row-Level Security -- RLS) cho Khía cạnh K2. Khi RLS được bật, chính bộ lập kế hoạch và thực thi truy vấn của CSDL sẽ tự động áp điều kiện lọc tenant\_\allowbreak{}id vào mọi thao tác đọc, thêm, sửa, xóa.

Để xác định phạm vi bảo vệ của RLS trên các bảng nghiệp vụ dùng chung, chính sách bảo mật dòng tuân theo bốn nguyên tắc sau:

\begin{enumerate}
\def\labelenumi{\arabic{enumi}.}
\item
  \textbf{Áp dụng cho chủ sở hữu bảng trong phạm vi PostgreSQL:} Thiết lập FORCE ROW LEVEL SECURITY trên ba bảng nghiệp vụ để chủ bảng cũng chịu chính sách khi truy cập thông thường \cite{v4-22}. Tách vai trò di trú khỏi vai trò ứng dụng có NOBYPASSRLS. Siêu người dùng và vai trò có BYPASSRLS vẫn có thể bỏ qua RLS.
\item
  \textbf{Kiểm tra dòng hiện có và giá trị mới:} Mệnh đề USING xác định các dòng hiện có mà thao tác được phép truy cập. WITH CHECK xác định giá trị mới được phép ghi khi chèn hoặc cập nhật \cite{v4-22}.
\item
  \textbf{Thiết kế theo nguyên tắc đóng an toàn mặc định (Fail-Closed):} Biểu thức kiểm tra (dùng giả mã nguyên lý trung tính GET\_\allowbreak{}TX\_\allowbreak{}CONTEXT để đọc biến phiên giao dịch) được viết sao cho nếu biến thuê bao chưa được gán hoặc rỗng, kết quả trả về là NULL thay vì báo lỗi hay bỏ qua lọc:
\end{enumerate}

\textbf{\emph{RLS\_\allowbreak{}Predicate: record.tenant\_\allowbreak{}id = NULLIF(GET\_\allowbreak{}TX\_\allowbreak{}CONTEXT(\textquotesingle tenant\_\allowbreak{}id\textquotesingle), \textquotesingle\textquotesingle)}}

Khi ngữ cảnh thuê bao là NULL, phép so sánh với record.tenant\_\allowbreak{}id cho kết quả UNKNOWN theo logic ba giá trị. Chính sách chỉ chấp nhận dòng khi biểu thức trả TRUE. Giá trị UNKNOWN do đó loại dòng khỏi tầm nhìn của thao tác hoặc làm kiểm tra giá trị mới thất bại, tùy mệnh đề và thao tác.

\begin{enumerate}
\def\labelenumi{\arabic{enumi}.}
\setcounter{enumi}{3}
\tightlist
\item
  \textbf{Đưa cột tenant\_\allowbreak{}id xuống tất cả các bảng con:} Thay vì chỉ lưu tenant\_\allowbreak{}id ở bảng cha rồi nối bảng JOIN khi kiểm tra RLS ở các bảng con, mọi bảng nghiệp vụ dùng chung đều lưu trực tiếp cột tenant\_\allowbreak{}id NOT NULL để rào chắn RLS kiểm tra trực tiếp trên dòng, giảm nhu cầu nối bảng. Hiệu năng truy vấn phụ thuộc vào kế hoạch thực thi và phân bố dữ liệu.
\end{enumerate}

Để đọc được mã thuê bao trong CSDL, biến phiên cần nhận giá trị từ TenantContext ở Tầng 8 truyền xuống (K1). Tuy nhiên, vì kết nối vật lý từ Bể kết nối dùng chung tồn tại lâu hơn một yêu cầu, nếu gán biến phiên ở mức kết nối vật lý thì giá trị tenant\_\allowbreak{}id cũ sẽ còn sót lại khi kết nối được cấp cho giao dịch khác.

Cách khắc phục là chỉ gắn biến phiên theo phạm vi một giao dịch duy nhất (ký hiệu giả mã SET\_\allowbreak{}TX\_\allowbreak{}CONTEXT trên Hình~\ref{fig:v4-3-9}, tương ứng lệnh SET LOCAL trong PostgreSQL \cite{v4-22}). Biến tenant\_\allowbreak{}id được gán ngay sau khi mở giao dịch (BEGIN) và tự động hủy khi giao dịch kết thúc (COMMIT hoặc ROLLBACK), trả kết nối sạch về Bể kết nối dùng chung như trên Hình~\ref{fig:v4-3-9}:

\VfourFigure{3.9}{assets/figures/v4/image25.png}{Tái sử dụng kết nối vật lý an toàn bằng biến phiên giao dịch.}{fig:v4-3-9}

Kết hợp FORCE RLS với biến cấu hình phạm vi giao dịch tạo thêm một lớp kiểm soát bên cạnh điều kiện lọc của ứng dụng. Việc tái sử dụng kết nối an toàn phụ thuộc vào thiết lập biến trong đúng giao dịch và việc kết thúc giao dịch trước khi trả kết nối về bể.

\subsection{Lưu trữ dữ liệu bán cấu trúc và tối ưu chỉ mục}\label{v4-sec-3-4-2}

Để nhiều thuê bao dùng chung một bảng vật lý cố định mà mỗi khách hàng vẫn tự thêm được các trường thông tin riêng (K4), các cách cũ như tạo sẵn hàng chục cột trống (col\_\allowbreak{}1..col\_\allowbreak{}20) hay mô hình EAV (Entity--Attribute--Value) đều gây tốn dung lượng hoặc phải nối quá nhiều bảng \cite{v4-21}. Thay vào đó, Tầng 6 áp dụng mô hình Lược đồ Lai Quan hệ -- Tài liệu: giữ các cột quan hệ chuẩn cho dữ liệu chung (bảng, cột, khóa ngoại) và thêm một cột kiểu tài liệu nhị phân (như JSONB) để chứa các trường mở rộng của từng tổ chức.

Để cột bán cấu trúc này chạy hiệu quả, thiết kế kết hợp hai việc: (1) đánh chỉ mục đảo ngược (như chỉ mục GIN) trực tiếp lên cột tài liệu để tìm kiếm nhanh theo cặp khóa -- giá trị bên trong; và (2) giao việc kiểm tra tên trường, kiểu dữ liệu và giá trị bắt buộc cho Động cơ Đánh giá Quy tắc ở Tầng 9 (Mục~\ref{v4-sec-3-3-2}) trước khi ghi vào CSDL.

Nhờ kết hợp cột quan hệ chuẩn với cột JSONB (kèm chỉ mục đảo ngược), hệ thống tránh được việc phải sửa bảng bằng ALTER TABLE gây khóa bảng, tránh lãng phí các cột dự phòng bỏ trống và tránh phải nối bảng phức tạp như mô hình EAV.

Với các truy vấn lọc theo thuê bao và sắp xếp theo một hoặc nhiều cột, chỉ mục tổ hợp có thể đặt tenant\_\allowbreak{}id ở đầu, ví dụ INDEX (tenant\_\allowbreak{}id, \textless các\_\allowbreak{}cột\_\allowbreak{}sắp\_\allowbreak{}xếp\textgreater). Cách đặt cột đầu hỗ trợ các mẫu truy vấn theo thuê bao. Các chỉ mục khác, như chỉ mục dùng để tìm theo tên miền phụ, phục vụ điều kiện khác.

Bộ lập kế hoạch có thể chọn quét dải chỉ mục cho truy vấn có điều kiện phù hợp. Lựa chọn còn phụ thuộc vào độ chọn lọc, kích thước bảng và ước lượng chi phí. Cần kế hoạch truy vấn và phép đo để đánh giá ảnh hưởng của chỉ mục và giới hạn thời gian chạy trong từng cấu hình.

\section{Phối hợp liên tầng và ánh xạ công nghệ}\label{v4-sec-3-5}

Chương~\ref{v4-chapter-3} đã kết nối bốn cụm tầng của mô hình Pool thành ba chốt chặn phối hợp nhau để bảo đảm trọn vẹn bốn khía cạnh kiểm soát K1--K4 (minh họa trên Hình~\ref{fig:v4-3-10}):

\begin{landscape}
\renewcommand{\VfourImageMaxHeight}{.55\textheight}
\begin{enumerate}
\def\labelenumi{\arabic{enumi}.}
\item
  Tại Tầng 7 (phần mềm trung gian): Xác thực JWT và đối chiếu ngữ cảnh thuê bao (K1), sau đó áp dụng hạn mức lưu lượng riêng ở lớp lọc ứng dụng (K3).
\item
  Tại Tầng 8--9 (ứng dụng và thực thi): Gắn ngữ cảnh theo phạm vi xử lý và thu hồi trong khối finally (K1). Bộ nhớ đệm dùng khóa phân vùng theo thuê bao (K2, K3), còn siêu dữ liệu cấu hình cho phép tùy biến ba lớp khi chạy (K4).
\item
  Tại Tầng 4--6 (cơ sở dữ liệu): Thiết lập ngữ cảnh theo giao dịch (K1). Chính sách RLS kiểm tra dòng trên các bảng nghiệp vụ (K2). Chỉ mục hỗ trợ một số mẫu truy vấn theo thuê bao (K3) và JSONB lưu các trường mở rộng (K4).
\end{enumerate}
\VfourFigure{3.10}{assets/figures/v4/image26.png}{Chuỗi phối hợp liên tầng của bốn khía cạnh kiểm soát K1-K4.}{fig:v4-3-10}

Để chuyển giao các nguyên lý trừu tượng K1-K4 sang giai đoạn cài đặt kỹ thuật ở Chương~\ref{v4-chapter-4} theo nhiều lựa chọn công nghệ khác nhau, Bảng~\ref{tab:v4-3-1} tổng hợp ánh xạ quy đổi sang ba tổ hợp công nghệ phổ biến:
\begingroup\VfourTableSetup
\renewcommand{\thetable}{3.1}
\begin{longtable}{@{}P{\dimexpr 0.25000\linewidth-7.500bp\relax}P{\dimexpr 0.25000\linewidth-7.500bp\relax}P{\dimexpr 0.25000\linewidth-7.500bp\relax}P{\dimexpr 0.25000\linewidth-7.500bp\relax}@{}}
\caption{Ánh xạ các nguyên lý kiểm soát K1-K4 sang một số tổ hợp công nghệ.}\label{tab:v4-3-1}\\
\toprule
Chốt chặn \& Nguyên lý & Java / Spring Boot + PostgreSQL (Ứng dụng) & Node.js (NestJS) + MySQL / PostgreSQL & Go + NoSQL (DocumentDB / MongoDB) \\
\midrule\endfirsthead
\toprule
Chốt chặn \& Nguyên lý & Java / Spring Boot + PostgreSQL (Ứng dụng) & Node.js (NestJS) + MySQL / PostgreSQL & Go + NoSQL (DocumentDB / MongoDB) \\
\midrule\endhead
\midrule\endfoot
\bottomrule\endlastfoot
K1 - Ngữ cảnh Luồng & JwtTenantFilter + ThreadLocal cố định trong khối try/finally & Middleware + AsyncLocalStorage (.run) tự động đóng phạm vi & Middleware + context.WithValue truyền qua context.Context \\
K2 - Cô lập Dữ liệu & FORCE RLS + SET LOCAL app.current\_\allowbreak{}tenant\_\allowbreak{}id trong giao dịch & PostgreSQL SET LOCAL hoặc View bảo mật / Repository Wrapper cưỡng chế tenant\_\allowbreak{}id & Logical tenant\_\allowbreak{}id Partition Key bắt buộc ở lớp DAO + Shard Key (tenant\_\allowbreak{}id) \\
K3 - Hiệu năng \& Cache & Redis Token Bucket (Lua) + Caffeine L1 (W-TinyLFU) + Composite B-Tree (tenant\_\allowbreak{}id, ...) & Redis Rate Limiter + In-memory LRU/LFU L1 + Chỉ mục phức hợp (tenant\_\allowbreak{}id, ...) & Redis/Envoy Token Bucket + Local Cache + Compound Index dẫn đầu tenant\_\allowbreak{}id \\
K4 - Tùy biến Động & Động cơ Rule 3 mức + cột JSONB (GIN Index) + Redis Pub/Sub đồng bộ L1 & Metadata JSON Ruleset + cột JSON / JSONB + Redis Pub/Sub & Document Dynamic Schema Validator + Native Flexible Fields + NATS/Redis Pub/Sub \\
\end{longtable}
\endgroup
\end{landscape}
